%%%%%%%%%%%%%%%%%%%%%%%%%%%%%%%%%%%%%%%%%%%%%%%%%%%%%%%%%%%%%%%%%%%%%%%%
%                                                                      %
%     File: Thesis_Glossary.tex                                        %
%     Tex Master: Thesis.tex                                           %
%                                                                      %
%     Author: Andre C. Marta                                           %
%     Last modified : 27 Feb 2024                                      %
%                                                                      %
%%%%%%%%%%%%%%%%%%%%%%%%%%%%%%%%%%%%%%%%%%%%%%%%%%%%%%%%%%%%%%%%%%%%%%%%
%
% The definitions can be placed anywhere in the document body
% and their order is sorted by <key> automatically when
% calling makeindex in the makefile
%
% ----------------------------------------------------------------------
% To create a glossary entry, use the following syntax:
%
% \newglossaryentry{<label>}{name={<key>}, description={<value>}}
%
% where the parameters are:
% <label> is the label of the entry,
% <key> is the acronym to be defined by the glossary entry (in lowercase, preferably)
% <value> is the actual definition of the current term
%
% To produce the desired term in the document, that will be replaced by
% the user-defined in the output, use the following syntax:
%
% \gls{<label>}
%
% ----------------------------------------------------------------------
% To create a acronym entry, use the following syntax:
%
% \newacronym{⟨label⟩}{⟨abbrv⟩}{⟨full⟩}
%
% where the parameters are:
% <label> is the label of the entry,
% <abbrv> is the acronym,
% <full> is the definition of the acronym
%
% To produce the desired term in the document, that will be replaced by
% the user-defined in the output, use one of the following syntaxes:
%
% \acrlong{<label>}
% \acrshort{<label>}
% \acrfull{<label>}
%
% ----------------------------------------------------------------------
% By default, only those entries defined in the main document using the
% commands above will be displayed in the glossary (list of acronyms),
% unless the command \glsaddall is used,
% ----------------------------------------------------------------------

% The order of the definitions below is irrelevant
% since the glossary is automatically ordered alphabetically

\newacronym{uav}{UAV}{Unmanned Aerial Vehicle}
\newacronym{usv}{USV}{Unmanned Surface Vessel}
\newacronym{mpc}{MPC}{Model Predictive Control}
\newacronym{nmpc}{NMPC}{Non-linear Model Predictive Control}
\newacronym{emsa}{EMSA}{European Maritime Safety Agency }
\newacronym{lti}{LTI}{linear time-invariant}
\newacronym{pid}{PID}{proportional–integral–derivative}
\newacronym{siso}{SISO}{single-input single-output}

\newacronym{lqr}{LQR}{linear quadratic regulator}
\newacronym{lqg}{LQG}{linear quadratic Gaussian}

\newacronym{ilc}{ILC}{iterative learning control}
\newacronym{rl}{RL}{reinforcement learning}


%\newacronym{⟨label⟩}{⟨abbrv⟩}{⟨full⟩}
%\newacronym{⟨label⟩}{⟨abbrv⟩}{⟨full⟩}

% ----------------------------------------------------------------------
% displays all entries (even those unused with commands \acrlong/short/full)
\glsaddall

% ----------------------------------------------------------------------
% vertical aligment of acronyms' long names
\setlength\LTleft{0pt}
\setlength\LTright{0pt}
\setlength\glsdescwidth{1.0\hsize}
